%% kapa: exam=1 round=35 year=2024 period=1 form=A subject=real_estate subject_ko=부동산학원론 seq=34 number=114
%% source: Q-Net 감정평가사 기출문제 (cst00309) artlSeq=5213229 file=raw/1008/5213229/1교시 문제지 원본.pdf
%% conversion: kordoc 4.24.51 -> tools/gen1.py (auto); review=proofread-sonnet; flags=none
\begin{question}{114}
부동산투자에 관한 설명으로 옳지 않은 것은? (단, 주어진 조건에 한함)
\begin{choices}
  \item 영업비용비율(OER)은 운영경비(OE)를 유효총소득(EGI)으로 나눈 비율이다.
  \item 총부채상환비율(DTI)이 높을수록 차입자의 부채상환가능성이 낮아진다.
  \item 채무불이행률(DR)은 유효총소득(EGI)으로 운영경비(OE)와 부채서비스(DS)를 감당할 수 있는 정도를 나타낸다.
  \item 총투자수익률(ROI)은 총투자액을 순영업소득(NOI)으로 나눈 비율이다.
  \item 지분투자수익률(ROE)은 세후현금흐름(ATCF)을 지분투자액으로 나눈 비율이다.
\end{choices}
\end{question}
\answer{4}
